\section{Measure work at the function boundaries}
A complete digest can be correct even when its implementation performs unnecessary
work. Count instructions only after establishing correctness, and retain the same
inputs and call boundary for each comparison. Here the counter starts at each function
entry and stops before its supplied return continuation. Assembly, session creation,
argument setup, caller stores and file I/O are excluded. The counts describe simulated
instruction work; they are not host elapsed time or a hardware cycle estimate.

The first implementation uses word rotates. The second expresses each rotate with two
word shifts and an OR, and uses alternative Boolean expressions for choice and majority.
Both use the same buffering, finalization and supplied return-frame scheme. They therefore produce the same
digests while executing different numbers of instructions in compression.

\begin{table}[htbp]
  \centering\caption{Executed instructions for update and finalization.}
  \tagpdfsetup{table/header-rows={1}}
  \begin{tabular}{lrrrr}
    \toprule
    Input      & Update & Update, shifts & Final & Final, shifts \\
    \midrule
    Empty      & 27     & 27             & 4393  & 5673          \\
    ASCII abc  & 60     & 60             & 4369  & 5649          \\
    55 bytes   & 632    & 632            & 3953  & 5233          \\
    56 bytes   & 643    & 643            & 8225  & 10785         \\
    64 bytes   & 4498   & 5778           & 4393  & 5673          \\
    65 bytes   & 4509   & 5789           & 4385  & 5665          \\
    4096 bytes & 286171 & 368091         & 4393  & 5673          \\
    \bottomrule
  \end{tabular}
\end{table}

Initialization takes 70 instructions in both versions. Schedule construction alone
takes 1318 instructions with rotates and 1702 with shifts and ORs. One compression call,
including its schedule construction, takes 3759 and 5039 instructions respectively.
The table's 55-, 56-, 64- and 65-byte inputs contain byte values equal to their indices;
the 4096-byte input uses byte value \((73i+128)\bmod256\). These measured inputs are
identified separately from other patterned correctness cases.

The finalization count roughly doubles between 55 and 56 bytes because the padding now
requires two compression blocks. It falls again at 64 bytes because update has already
compressed the message block, leaving only one final padding block. Empty-message
finalization also performs a compression. These jumps follow block structure, rather
than a smooth cost per message byte.

For an optimization experiment, first choose which count is supposed to change.
Avoiding the copy of an already complete source block targets update. A circular
schedule targets workspace size and may change schedule access costs. Unrolling the
round loop targets loop overhead but increases code size. Test each change against
unaligned input, split messages, wide pointers and memory ownership before interpreting
its work count. Equivalent implementations are not required to reproduce these exact
instruction totals.
